\section{本章小结}

本章介绍了随机变量的数字特征：
\begin{itemize}
    \item 对于一个变量，有：数学期望、方差，描述了变量的分布情况；
    \item 对于两个变量，有：协方差、相关系数，描述了样本点$\left\{ X=x_i,Y=y_j \right\} $的分布是否趋于直线。
\end{itemize}

~

数学期望定义公式
\begin{align*}
&EX:=\sum_i{x_i\cdot p_i} \\
&EX:=\int_{-\infty}^{+\infty}{xf\left( x \right) \cdot dx}
\end{align*}

方差定义公式和便捷公式
\begin{align*}
&DX:=E\left[ \left( X-EX \right) ^2 \right] \\
&DX=E\left( X^2 \right) -\left( EX \right) ^2
\end{align*}

协方差定义公式和便捷公式
\begin{align*}
&\mathrm{Cov}\left( X,Y \right) :=E\left[ \left( X-EX \right) \left( Y-EY \right) \right] \\
&\mathrm{Cov}\left( X,Y \right) =E\left( XY \right) -EX\cdot EY
\end{align*}

相关系数定义公式和便捷公式
\begin{align*}
&\rho _{XY}:=\frac{\mathrm{Cov}\left( X,Y \right)}{\sqrt{DX}\cdot \sqrt{DY}} \\
&\rho _{XY}=\frac{E\left( XY \right) -EX\cdot EY}{\sqrt{E\left( X^2 \right) -\left( EX \right) ^2}\cdot \sqrt{E\left( Y^2 \right) -\left( EY \right) ^2}}
\end{align*}

数学期望的性质
\begin{align*}
&E\left( c \right) =c \\
&X\geqslant a\Rightarrow EX\geqslant a,X\leqslant a\Rightarrow EX\leqslant a \\
&E\left( kX \right) =k\cdot E\left( X \right) \\
&E\left( kX+c \right) =k\cdot E\left( X \right) +c \\
&E\left( X\pm Y \right) =E\left( X \right) \pm E\left( Y \right) \\
&\text{若}X,Y\text{相互独立}\Rightarrow E\left( XY \right) =E\left( X \right) \cdot E\left( Y \right) \\
&\left[ E\left( XY \right) \right] ^2\leqslant E\left( X^2 \right) \cdot E\left( Y^2 \right)
\end{align*}

方差的性质
\begin{align*}
&D\left( c \right) =0 \\
&D\left( X \right) \geqslant 0, \text{当且仅当}P\left\{ X=EX \right\} =1\text{时等号成立} \\
&D\left( kX \right) =k^2\cdot D\left( X \right) \\
&D\left( kX+c \right) =k^2\cdot D\left( X \right) \\
&\text{若}X,Y\text{相互独立}\Rightarrow D\left( X\pm Y \right) =D\left( X \right) +D\left( Y \right)
\end{align*}

协方差的性质
\begin{align*}
&\mathrm{Cov}\left( X,X \right) =DX \\
&D\left( X\pm Y \right) =DX+DY\pm 2\mathrm{Cov}\left( X,Y \right) \\
&\mathrm{Cov}\left( X,c \right) =0 \\
&\mathrm{Cov}\left( X,Y \right) =\mathrm{Cov}\left( Y,X \right) \\
&\mathrm{Cov}\left( aX,bY \right) =ab\mathrm{Cov}\left( X,Y \right) \\
&\mathrm{Cov}\left( X_1\pm X_2,Y \right) =\mathrm{Cov}\left( X_1,Y \right) \pm \mathrm{Cov}\left( X_2,Y \right) \\
&\left[ \mathrm{Cov}\left( X,Y \right) \right] ^2\leqslant DX\cdot DY
\end{align*}

相关系数的性质
\begin{align*}
&\rho _{XY}\in \left[ -1,1 \right] \\
&X\text{和}Y\text{成线性关系}\Leftrightarrow \rho _{XY}=\pm 1 \\
&\rho _{\left( aX \right) \left( bY \right)}=\begin{cases}
	\rho _{XY},&		ab>0\\
	-\rho _{XY},&		ab<0\\
\end{cases}
\end{align*}

~

数字特征的物理意义和线性代数意义：
\begin{table}[h]
\centering
\begin{tabular}{lll}
    \toprule
    数字特征 & 物理意义/线性代数意义 & 量纲\\
    \midrule
    数学期望 & 质心位置 & $\left[ \mathrm{m} \right] $\\
    方差 & 单位质量的转动惯量，范数 & $\left[ \mathrm{m}^2 \right] $\\
    协方差 & 两个向量的内积 & $\left[ \mathrm{m}^2 \right] $\\
    相关系数 & 两个向量的夹角的余弦 & 无量纲\\
    \bottomrule
\end{tabular}
\end{table}

学习概率论中的数字特征，要特别注意它们的物理意义和几何意义，用其他学科的知识点相互印证，加深对这些概念的理解。




